\section{Rank-one ports with arbitrary coordinate widths}\label{sec:rank-one}
The reversible-rule bounds require an input to be recoverable from the
other symbols in its row. A singular coefficient need not permit this.
We now allow arbitrary finite coordinate widths, provided every recovery
block has rank at most one. Throughout this section, $K$ is any field and
$G=(V,E)$ is a finite simple triangle-free graph of maximum degree three.

\begin{theorem}[R: intrinsic half bound]\label{thm:rank-one}
Let $A_v$ be finite-dimensional $K$-vector spaces and let
$C\leq\prod_{v\in V}A_v$ be a linear code. Suppose linear maps
$T_{vu}:A_u\to A_v$ satisfy
\[
 x_v=\sum_{u\in V}T_{vu}x_u\quad(x\in C),\qquad
 T_{vu}=0\text{ if }vu\notin E,\qquad \rank T_{vu}\leq1.
\]
Writing $\pi_v$ for coordinate projection, we have
\[
 2\dim C\leq\sum_{v\in V}\dim\pi_v(C)\leq\sum_{v\in V}\dim A_v.
\]
Consequently, if $B$ is the block operator with blocks $T_{vu}$ and
$N=\sum_v\dim A_v$, then $2\rank(I-B)\geq N$.
\formal{RankOneCodes.lean}{LowRank.rank_one_code_half,LowRank.rank_one_code_capacity,LowRank.rank_one_block_rank,LowRank.scalar_recovery_half}
\end{theorem}
\begin{outline}
The proof first establishes a statement about an arrangement of subspaces
$U_v\leq W$ spanning $W$. Its supplied component maps
$R_{vu}:U_v\to W$ have images in $U_u$, vanish off $G$, have rank at most
one, and sum to the inclusion of $U_v$ in $W$.
It proves $2\dim W\leq\sum_v\dim U_v$.

Induct on $\dim W$. Quotienting by any nonzero subspace $S$ produces
another arrangement of the same kind: choose a linear section of each
local quotient image and compose each old component with that section
and the ambient quotient. This constructs new maps of no larger rank;
it does not assume that the old maps descend. A counterexample would
therefore satisfy
\[
 f(S):=2\dim S-\sum_v\dim(U_v\cap S)\geq1\qquad(S\ne0).
\]
Distinct coordinate spaces are disjoint. Each coordinate has dimension
at most two: its at most three image lines span it, and equality with a
three-dimensional row span would contradict that disjointness.

The submodularity of $f$ makes intersecting spaces of deficit one have
deficit-one intersection and sum. Equality also gives the exact identity
$U_v\cap(S+T)=(U_v\cap S)+(U_v\cap T)$ for those spaces.
A plane coordinate's three input lines span a three-dimensional space
of deficit one, meeting no other coordinates. Two adjacent plane rows
would have a tight intersection supported on just two lines, an
impossible odd deficit. Thus every neighbor of a plane is a line.

Starting with one plane's row span, grow a tight space by either adjoining
two outside lines whose nonzero row contributions sum into the space,
or absorbing a plane's row span along one shared input line. Exact
intersection accounting preserves whole-coordinate membership and the
inclusion of every known plane's row span. Every known line retains a
known neighbor. A reachable space of maximum dimension is closed under
all the supplied component maps: a partially met plane span would admit
an absorption, and a known line has at most two outside neighbors, whose
nonzero contributions would admit a line extension.

If this closed space is proper, its coordinate intersections span it
and the actual recovery maps restrict to it. Its deficit one contradicts
induction. If it is the whole space, inspect the last growth operation.
An initial star or a final line extension contradicts recovery of one
of its leaves. For a final star absorption, the two newly added input
axes are independent modulo the predecessor space. Triangle-freeness
and this independence make that predecessor closed under all the old
rows, again contradicting induction. If there was no plane to begin
with, evaluation on line generators gives a scalar code; the scalar
half bound applies, counting only nonzero coordinates.

For the stated code, set $W=C^*$ and $U_v=\operatorname{im}\pi_v^*$.
Local linear sections of $\pi_v^*$ turn the dual recovery equations
into the required component maps without increasing their ranks.
These coordinate duals span $C^*$, and $\dim U_v=\dim\pi_v(C)$.
The ambient-capacity bound follows by inclusion, and rank--nullity gives
the block-operator statement. Width one recovers the scalar graph
half bound.
\end{outline}

The arrangement theorem and its global closed-subspace step are
\formal{RankOneBound.lean}{LowRank.rank_one_arrangement_half,LowRank.Core.exists_smaller_closed}
The generic quotient, restriction, and code-dual constructions are shared
in the repository's \texttt{lean-machine} library. In particular, no
compatible choice of a scalar basis across different ports is assumed.
For the construction program, a code exceeding half capacity must escape
the rank-one condition somewhere. This is a restriction on the recovery
maps; it is not a condition supplied by tape geometry.

\begin{theorem}[Paid exceptional rows]\label{thm:rank-one-exceptions}
Consider a spanning arrangement $U_v\leq W$ with supplied recoveries as
above, but allow unrestricted component ranks at vertices in $D\subseteq V$.
Every row outside $D$ still has rank-at-most-one components. Set
$H=\sum_{v\in D}U_v$ and $\Delta=2\dim W-\sum_v\dim U_v$. Then
\[
 \Delta\leq2\dim H-\sum_v\dim(U_v\cap H)
 \leq\dim H\leq\sum_{v\in D}\dim U_v.
\]
\formal{RankOneCodes.lean}{LowRank.rank_one_exception_defect,LowRank.exceptional_span_defect}
\end{theorem}
\begin{outline}
Quotient by $H$, using the local-section construction. Exceptional
coordinates become zero, and all surviving components have rank at
most one. Apply the arrangement half bound to this quotient. The
identities $\dim(U_v\bmod H)=\dim U_v-\dim(U_v\cap H)$ give the first
inequality. The sum of the exceptional coordinate dimensions is at
least $\dim H$ and contributes in full to the intersection sum, giving
the remaining inequalities.
\end{outline}
Thus any excess above half must be paid by the actual span of exceptional
coordinates. The theorem does not assert that this span is small in a
native tape computation.
